% !TEX program = xelatex
\documentclass{misisarticle}

\grnti{20.23.17, 27.41.41}
\asjc{1710, 1703}
\oecd{1020}

\title{Поиск нечётких дубликатов в клиентской базе данных на основе расстояния Левенштейна}

\misisauthor{Е.~А.~Семенов}{кафедра инженерной кибернетики НИТУ «МИСиС»}{Москва, Россия}{m1905644@edu.misis.ru}

\abstract{%
Рассматривается задача поиска дублирующихся карточек клиентов в клиентской базе данных коммерческой организации. Дубликаты возникают из"=за ошибок ручного ввода (опечаток, ошибок регистра, неверной раскладки) и не выявляются стандартными средствами СУБД, ориентированными на точное совпадение. Выполнен обзор методов нечёткого сравнения строк и подходов к связыванию записей, обоснован выбор расстояния Левенштейна для полей «Фамилия», «Имя», «Отчество». Предложен критерий отнесения записи к дубликатам по порогу схожести и реализована программа поиска на языке Python. Показано, что критерий находит записи с одиночными опечатками, но ошибочно считает дубликатами мужские и женские формы ФИО; предложены пути устранения этого недостатка.
}

\keywords{%
дедупликация, нечёткий поиск, расстояние Левенштейна, клиентская база данных, связывание записей
}

\begin{document}
\maketitle

\section{Введение}

Клиентская база данных (БД) является ключевым информационным ресурсом коммерческой организации: на её основе ведутся продажи, начисляются бонусы по программам лояльности, выполняются рассылки и анализ аудитории. Карточки клиентов создаются вручную несколькими операторами, а также поступают через сайт и при импорте из внешних источников. При таком ведении базы один клиент нередко регистрируется несколько раз: оператор не находит существующую карточку из"=за опечатки, ошибки регистра или замены кириллических букв латинскими~\cite{dyadichev2009}.

Дубликаты приводят к повторным рассылкам и звонкам, разделению истории покупок и бонусов между карточками, завышению численности клиентской базы. Кроме того, закон «О персональных данных» требует от оператора обеспечивать точность обрабатываемых данных~\cite{fz152}. Стандартные СУБД и механизмы учётных систем ориентированы на поиск точных совпадений, поэтому для выявления таких записей необходим нечёткий поиск.

Цель работы "--- разработать метод и программу поиска дублирующихся записей о клиентах на основе расстояния Левенштейна. Объект исследования "--- клиентская БД, предмет "--- алгоритмы поиска нечётких дубликатов на основе редакционного расстояния.

\section{Анализ существующих работ}

\subsection{Дедупликация и ошибки ручного ввода}

Под дедупликацией в системах хранения данных понимают поиск повторяющихся данных и замену их указателем либо удаление~\cite{lyashchenko2013}. Такие механизмы работают на уровне файлов и блоков и выявляют только побитовые совпадения, поэтому карточки клиентов с различающимся содержимым полей они не обнаруживают.

Распространённость ошибок ручного ввода подтверждается исследованием поисковых запросов~\cite{yandex}: чаще всего встречаются замены «и"--"е» (27\,\%) и «а"--"о» (25\,\%), лишние или пропущенные пробелы, потеря и удвоение букв. Все эти ошибки сводятся к вставке, удалению или замене одного символа, что делает естественным применение редакционного расстояния.

\subsection{Методы нечёткого сравнения строк}

Методы нечёткого сравнения можно разделить на группы (табл.~\ref{tab:methods}). Коэффициент Жаккара оценивает долю общих символов и не учитывает их порядок. Редакционные метрики (Левенштейна, Дамерау"--"Левенштейна, Джаро"--"Винклера) подсчитывают число операций редактирования и хорошо выявляют опечатки. N-граммные методы и алгоритм шинглов применяются в основном для длинных текстов~\cite{williams2013, zhang2010}. Фонетические алгоритмы (Soundex, Metaphone) не предназначены для выявления случайных опечаток, а методы машинного обучения требуют моделей и значительных вычислительных ресурсов, что для коротких полей избыточно.

\begin{table}[htbp]
\caption{Классификация методов поиска похожих строк}
\label{tab:methods}
\centering
\small
\begin{tabular}{p{2.4cm}p{2.6cm}p{2.5cm}}
\toprule
\textbf{Группа} & \textbf{Представители} & \textbf{Особенность} \\
\toprule
Общие символы & коэф. Жаккара & не учитывает порядок \\
Редакционное расстояние & Левенштейн, Джаро"--"Винклер & хорошо находит опечатки \\
N-граммы & шинглы & длинные тексты \\
Фонетические & Soundex, Metaphone & не для опечаток \\
Машинное обучение & векторные представления & требует ресурсов \\
\bottomrule
\end{tabular}
\end{table}

\subsection{Связывание записей}

Задача поиска записей об одном объекте реального мира известна как связывание записей (record linkage). Вероятностная постановка предложена Феллеги и Сантером~\cite{fellegi1969}: для пары записей вычисляется вектор сходства по полям, и пара относится к совпадениям, несовпадениям или спорным случаям, требующим ручной проверки. В монографии~\cite{christen2012} описан полный цикл: нормализация данных, блокировка, сравнение полей, классификация пар и оценка качества. Вопросы быстрого приближённого сопоставления строк рассмотрены в~\cite{baeza1999}.

\section{Постановка задачи}

Необходимо по выбранной записи клиента найти в БД записи"=дубликаты, сравнивая поля «Фамилия», «Имя», «Отчество». Запись $r$ считается дубликатом искомой записи $q$, если схожесть по каждому из трёх полей превышает порог $T$:
\begin{equation}
\min_{k \in \{1,2,3\}} s(r_k, q_k) > T,
\label{eq:criterion}
\end{equation}
где $r_k$, $q_k$ "--- значения $k$-го поля, $s$ "--- функция схожести, выраженная в процентах. Дополнительные поля (дата рождения, телефон, e-mail) в данной постановке не используются.

\section{Описание метода}

\subsection{Расстояние Левенштейна}

Расстояние Левенштейна между строками $S_1$ и $S_2$ длиной $M$ и $N$ "--- минимальное число вставок, удалений и замен символа, переводящих одну строку в другую. Оно вычисляется заполнением матрицы $D$ размером $(M{+}1){\times}(N{+}1)$ по рекуррентной формуле
\begin{equation}
D(i,j)=\min\begin{cases}
D(i,j-1)+1,\\
D(i-1,j)+1,\\
D(i-1,j-1)+m(S_1[i],S_2[j]),
\end{cases}
\label{eq:lev}
\end{equation}
где $m(a,b)=0$ при $a=b$ и $1$ иначе; $D(i,0)=i$, $D(0,j)=j$. Искомое расстояние равно $D(M,N)$, сложность вычисления "--- $O(M \cdot N)$, что для коротких полей несущественно.

Для удобства пользователя результат выражают в процентах:
\begin{equation}
\mathrm{ratio}(S_1,S_2)=\frac{|S_1|+|S_2|-d^{*}(S_1,S_2)}{|S_1|+|S_2|}\cdot 100\,\%,
\label{eq:ratio}
\end{equation}
где $d^{*}$ "--- расстояние, в котором замена считается за две операции. Например, для «Авдеев» и «Авдев» $d^{*}=1$ и схожесть равна $10/11\approx 91\,\%$.

Метод выбран потому, что поля ФИО "--- короткие строки с локальными искажениями, он не требует обучающей выборки, а результат в процентах позволяет задать понятный порог. Для устранения влияния регистра и порядка слов используется функция \texttt{token\_sort\_ratio} библиотеки fuzzywuzzy.

\subsection{Программная реализация}

Программа написана на Python: pandas хранит БД в объекте DataFrame, tkinter обеспечивает графический интерфейс, fuzzywuzzy с python-Levenshtein выполняет сравнение строк. БД загружается из файла CSV. Пользователь выбирает запись, задаёт порог (по умолчанию 70\,\%) и запускает поиск; найденные дубликаты выводятся в отдельную таблицу, из которой их можно удалить. Для одной искомой записи выполняется не более $3n$ сравнений ($n$ "--- число записей), поэтому полный поиск всех пар дубликатов требует порядка $n^2$ сравнений, и при росте БД необходима оптимизация, например блокировка~\cite{christen2012}.

\subsection{Выбор порога и ограничения критерия}

По формуле~\eqref{eq:ratio} одна замена в имени из четырёх букв даёт 75\,\% («Иван"--"Ивон»), пропуск буквы в фамилии "--- 91\,\%, а разные имена "--- существенно меньше («Иван"--"Илья» "--- 25\,\%). Поэтому порог 70\,\% пропускает одиночные опечатки в каждом поле. На тестовой БД запись «Авдеев Алексей Ильич» встречается четыре раза, в трёх из них с ошибками (регистр, пропущенные и искажённые буквы); при пороге 70\,\% программа находит все четыре записи.

Вместе с тем критерий~\eqref{eq:criterion} не различает мужские и женские формы ФИО. Для пары «Петров Александр Сергеевич"~--- Петрова Александра Сергеевна» схожесть полей составляет 92, 95 и 78\,\%, поэтому брат и сестра (или супруги) будут ошибочно отнесены к дубликатам. Другие ограничения: отсутствие нормализации (латиница вместо кириллицы, «ё»$\to$«е»), использование только полей ФИО.

\section{Заключение}

Предложен и реализован метод поиска нечётких дубликатов в клиентской БД на основе расстояния Левенштейна. Он находит записи с опечатками и ошибками регистра в полях ФИО, но даёт ложные срабатывания на мужские и женские формы. Для их устранения необходимо учитывать уточняющие поля (дату рождения, телефон, e-mail), нормализовать значения перед сравнением и обосновать порог на размеченных данных; для сокращения числа сравнений на больших базах целесообразно применять блокировку.

\begin{thebibliography}{9}

\bibitem{fz152} Федеральный закон от 27.07.2006 № 152-ФЗ «О персональных данных». \textit{Собрание законодательства РФ}, 2006, № 31 (ч.~1), ст.~3451.

\bibitem{dyadichev2009} Дядичев В.\,В., Локотош Б.\,Н., Колесников А.\,В. \textit{Теория информации}: учебное пособие. Луганск: Изд-во ВНУ им. В.~Даля, 2009. 320 с.

\bibitem{lyashchenko2013} Лященко А.\,О. Використання систем дедуплікації даних. \textit{Технологический аудит и резервы производства}, 2013, т.~3, № 1(11).

\bibitem{yandex} Исследования Яндекс [Электронный ресурс]. URL: yandex.ru/company/researches/2012/ya\_orfo (дата обращения: 15.10.2026).

\bibitem{baeza1999} Baeza-Yates R., Navarro G. Faster approximate string matching. \textit{Algorithmica}, vol.~23, pp.~127--158, 1999.

\bibitem{williams2013} Williams K., Giles C.\,L. Near duplicate detection in an academic digital library. In \textit{Proc. ACM Symposium on Document Engineering}, pp.~91--94, 2013.

\bibitem{zhang2010} Zhang Q., Zhang Y., Yu H., Huang X. Efficient partial-duplicate detection based on sequence matching. In \textit{Proc. 33rd Int. ACM SIGIR Conf.}, pp.~675--682, 2010.

\bibitem{fellegi1969} Fellegi I.\,P., Sunter A.\,B. A theory for record linkage. \textit{J. American Statistical Association}, vol.~64, no.~328, pp.~1183--1210, 1969.

\bibitem{christen2012} Christen P. \textit{Data Matching: Concepts and Techniques for Record Linkage, Entity Resolution, and Duplicate Detection}. Berlin, Heidelberg: Springer, 2012. 270 p.

\end{thebibliography}

\end{document}

